\subsection{Artin's theorem}

THEOREM 2 (Artin). --- Let N be a field, \(\Gamma\) a group of automorphisms of N and K the field of invariants of \(\Gamma\) . Let V be a vector sub-K-space of N of finite dimension over K. Then every K-linear mapping u of V into N is a linear combination with coefficients in N of the restrictions to V of elements of \(\Gamma\) .

Let \(u\) be a K-linear mapping of \(\mathbf{V}\) into \(\mathbf{N}\) and let \(\mathrm{V}_{(\mathrm{N})} = \mathrm{N} \otimes_{\mathrm{K}} \mathrm{V}\) be the vector N-space derived from \(\mathbf{V}\) by extension of scalars; denote by \(\tilde{u}\) the N-linear form on \(\mathrm{V}_{(\mathrm{N})}\) such that \(\tilde{u}(x \otimes y) = x \cdot u(y)\) for \(x \in \mathbb{N}\) and \(y \in \mathbb{V}\) . For each \(\sigma \in \Gamma\) there exists an N-linear form \(h_{\sigma}\) on \(\mathrm{V}_{(\mathrm{N})}\) such that \(h_{\sigma}(x \otimes y) = x \sigma(y)\) for \(x \in \mathbb{N}\) and \(y \in \mathbb{V}\) . The canonical mapping of \(\mathrm{V}_{(\mathrm{N})} = \mathrm{N} \otimes_{\mathrm{K}} \mathrm{V}\) into \(\mathrm{N} \otimes_{\mathrm{K}} \mathrm{N}\) is injective. Now the Cor. of Prop. 8 (V, p.~64) show that the intersection of the kernels of the linear forms \(h_{\sigma}\) on \(\mathrm{V}_{(\mathrm{N})}\) is reduced to 0. Therefore (II, p.~302, Cor. 1) there exist \(\sigma_1, \ldots, \sigma_n\) in \(\Gamma\) and \(a_1, \ldots, a_n\) in \(\mathbb{N}\) such that \(\tilde{u} = \sum_{i=1}^{n} a_i h_{\sigma_i}\) whence \(u(x) = \sum_{i=1}^{n} a_i \sigma_i(x)\) for all \(x \in V\) .

Let us equip the set \(N^{N}\) of all mappings of N into N with the product topology of the discrete topologies of the factors. Th. 2 means that the set of linear combinations with coefficients in N of elements of \(\Gamma\) is dense in the set of K-linear mappings of N into itself.

THEOREM 3. --- Let N be a field, \(\Gamma\) a finite group of automorphisms of N and K the field of invariants of \(\Gamma\) . Let \(n\) be the cardinal of \(\Gamma\) .

\begin{enumerate}
\def\labelenumi{\alph{enumi})}
\item
  We have \([\mathbf{N}:\mathbf{K}] = n\) and \(\mathbf{N}\) is a Galois extension of \(\mathbf{K}\) with Galois group \(\Gamma\) .
\item
  Let \(\sigma_1, \ldots, \sigma_n\) be the elements of \(\Gamma\) and \((x_1, \ldots, x_n)\) a basis of \(\mathbf{N}\) over \(\mathbf{K}\) , then \(\det (\sigma_i(x_j)) \neq 0\) .
\item
  Let \(u\) be a K-linear mapping of \(\mathbf{N}\) into \(\mathbf{N}\) . There exists a unique family \((a_{\sigma})_{\sigma \in \Gamma}\) of elements of \(\mathbf{N}\) such that \(u(x) = \sum_{\sigma \in \Gamma} a_{\sigma} \sigma(x)\) for all \(x \in \mathbf{N}\) .
\end{enumerate}

We equip the ring \(N \otimes_{K} N\) with the N-algebra structure whose external law is given by \(\lambda(x \otimes y) = x \otimes \lambda y\) for \(\lambda, x, y\) in N. Then the dimension of the vector N-space \(N \otimes_{K} N\) is {[}N: K{]}. The dimension of the product vector N-space \(N^{\Gamma}\) is equal to n.~The mapping \(\psi\) defined in the Cor. of Prop. 8 (V, p.~64) is an N-isomorphism of \(N \otimes_{K} N\) onto \(N^{\Gamma}\) , whence {[}N: K{]} = n.~Let \(\Delta\) be the group of K-automorphisms of N. We have \(\Gamma \subset \Delta\) , hence K is the field of invariants of \(\Delta\) , and N is a Galois extension of K. Further, the order of \(\Delta\) is at most equal to {[}N: K{]} by Dedekind's theorem (V, p.~27, Cor. 2) and since the order of \(\Gamma\) equals {[}N: K{]}, we have \(\Gamma = \Delta\) . Hence \(\Gamma\) is the Galois group of N over K, and this proves a).

With the notation of \(b)\) put \(f_{i} = \psi(x_{i} \otimes 1)\) ; we have \(f_{i}(\sigma) = \sigma(x_{i})\) for \(1 \leqslant i \leqslant n\) and \(\sigma \in \Gamma\) . Since \(\psi\) is an isomorphism of vector N-spaces, the sequence \((f_{1}, \ldots, f_{n})\) is a basis of \(\mathbf{N}^{\Gamma}\) over \(\mathbf{N}\) , whence \(\det(f_{j}(\sigma_{i})) \neq 0\) , that is,

\[
\det \left(\sigma_ {i} (x _ {j})\right) \neq 0.
\]

This proves \(b\) ).

Finally, \(c)\) follows from Th. 2 (V, p.~65) which proves the existence of a family \((a_{\sigma})_{\sigma \in \Gamma}\) such that \(u(x) = \sum_{\sigma \in \Gamma} a_{\sigma} \sigma(x)\) (for all \(x \in \mathbf{N}\) ) and from Dedekind's theorem

(V, p.~27, Cor. 2) which proves the uniqueness of \((a_{\sigma})_{\sigma \in \Gamma}\) .

